\section{Рабочий проект}

\subsection{Модуль index.html}

Главная страница приложения. Содержит разметку всех функциональных разделов системы: авторизация, профиль пользователя, проверка документов, генерация титульных листов, статистика, отправка документов в РИО и административная панель. Выполняет подключение таблицы стилей и JavaScript-модуля.

\subsection{Модуль styles.css}

Таблица стилей интерфейса. Определяет оформление пользовательского интерфейса: цветовую схему, шрифты, размеры элементов, анимации, адаптивную вёрстку, стили прогресс-баров, индикаторов загрузки и модальных окон.

\subsection{Модуль app.js}

Основной сценарий app.js реализует бизнес-логику веб-приложения: обработку пользовательских событий, выполнение HTTP-запросов к REST API, управление токенами авторизации, динамическое обновление содержимого страниц и обработку результатов проверки документов. Ниже приведено описание ключевых функций, сгруппированных по функциональным блокам.

\subsubsection{Функции авторизации и управления состоянием}

В таблице~\ref{tab:js_auth} представлены функции авторизации.

\begin{xltabular}{\textwidth}{|l|l|X|}
	\caption{Функции авторизации\label{tab:js_auth}}\\
	\hline
	\textbf{Название} & \textbf{Доступ} & \textbf{Описание} \\
	\hline
	\endfirsthead
	
	\continuecaption{Продолжение таблицы \ref{tab:js_auth}}
	\hline
	\textbf{Название} & \textbf{Доступ} & \textbf{Описание} \\
	\hline
	\endhead
	
	\hline
	\endlastfoot
	
	handleLogin & public & Отправка учётных данных на сервер, сохранение JWT-токенов, переход в профиль. \\
	\hline
	handleRegister & public & Отправка данных регистрации, отображение сообщения об успехе. \\
	\hline
	logout & public & Отзыв refresh-токена на сервере, очистка локального состояния. \\
	\hline
	refreshAccessToken & internal & Автоматическое обновление access-токена по refresh-токену. \\
	\hline
	apiFetch & internal & Унифицированная отправка HTTP-запросов с обработкой авторизации. \\
	\hline
\end{xltabular}

\subsubsection{Функции проверки документов}

В таблице~\ref{tab:js_check} приведены функции проверки документов.

\begin{xltabular}{\textwidth}{|l|l|X|}
	\caption{Функции проверки документов\label{tab:js_check}}\\
	\hline
	\textbf{Название} & \textbf{Доступ} & \textbf{Описание} \\
	\hline
	\endfirsthead
	
	\continuecaption{Продолжение таблицы \ref{tab:js_check}}
	\hline
	\textbf{Название} & \textbf{Доступ} & \textbf{Описание} \\
	\hline
	\endhead
	
	\hline
	\endlastfoot
	
	initCheckSubnav & public & Инициализация переключения между типами документов (методические указания, учебное пособие, монография). \\
	\hline
	updateCheckUIType & internal & Изменение текстов подсказок и кнопок в зависимости от выбранного типа документа. \\
	\hline
	handleCheck & public & Загрузка файла, отправка POST-запроса на /checker/check, отображение результатов (количество ошибок, ссылки на скачивание). \\
	\hline
	showCheckLoading & internal & Отображение анимированного индикатора загрузки и прогресс-бара. \\
	\hline
	initFileUpload & public & Настройка кастомной области загрузки файлов с поддержкой drag-and-drop. \\
	\hline
\end{xltabular}

\subsubsection{Функции генерации титульных листов}

В таблице~\ref{tab:js_title} описаны функции генерации титульных листов.

	\begin{xltabular}{\textwidth}{|p{5.1cm}|p{1.6cm}|X|}
		\caption{Функции генерации титульных листов\label{tab:js_title}}\\
		\hline
		\textbf{Название} & \textbf{Доступ} & \textbf{Описание} \\
		\hline
		\endfirsthead
		
		\continuecaption{Продолжение таблицы \ref{tab:js_title}}
		\hline
		\textbf{Название} & \textbf{Доступ} & \textbf{Описание} \\
		\hline
		\endhead
		
		\hline
		\endlastfoot
		
		handleTitlePageGenerate & public & Сбор данных из формы, отправка POST-запроса на соответствующий эндпоинт (/title-page/generate, /generate-tutorial, /generate-monograph) и скачивание полученного DOCX-файла. \\
		\hline
		addManualReviewer & public & Добавление блока рецензента в форму методических указаний. \\
		\hline
		addTutorialReviewer & public & Добавление блока рецензента в форму учебного пособия. \\
		\hline
		addTutorialDirection & public & Добавление направления подготовки в форму учебного пособия. \\
		\hline
		addMonographAuthor & public & Добавление блока автора в форму монографии. \\
		\hline
		removeReviewer & public & Удаление блока рецензента (учебное пособие) с проверкой минимального количества. \\
		\hline
		removeManualReviewer & public & Удаление блока рецензента (методические указания). \\
		\hline
		removeMonographAuthor & public & Удаление блока автора (монография). \\
		\hline
	\end{xltabular}

\subsubsection{Функции отправки документов в РИО}

В таблице~\ref{tab:js_rio} представлены функции отправки документов в РИО.

\begin{xltabular}{\textwidth}{|l|l|X|}
	\caption{Функции отправки в РИО\label{tab:js_rio}}\\
	\hline
	\textbf{Название} & \textbf{Доступ} & \textbf{Описание} \\
	\hline
	\endfirsthead
	
	\continuecaption{Продолжение таблицы \ref{tab:js_rio}}
	\hline
	\textbf{Название} & \textbf{Доступ} & \textbf{Описание} \\
	\hline
	\endhead
	
	\hline
	\endlastfoot
	
	handleRioStep1 & public & Загрузка двух вспомогательных файлов (шаг 1), переход к шагу 2 при успехе. \\
	\hline
	handleRioStep2 & public & Загрузка основного DOCX-файла, запуск проверки, переход к шагу 3 при отсутствии ошибок. \\
	\hline
	handleRioStep3 & public & Отправка комментария и всех трёх файлов на email РИО, завершение процесса. \\
	\hline
	goToRioStep & internal & Переключение между шагами мастера, обновление индикаторов. \\
	\hline
	resetRioSession & public & Сброс сессии отправки (очистка временных файлов на сервере). \\
	\hline
\end{xltabular}

\subsubsection{Функции модулей статистики и администрирования}

В таблице~\ref{tab:js_stat_admin} приведены функции модулей статистики и администрирования.

\begin{xltabular}{\textwidth}{|l|l|X|}
	\caption{Функции модулей статистики и администрирования\label{tab:js_stat_admin}}\\
	\hline
	\textbf{Название} & \textbf{Доступ} & \textbf{Описание} \\
	\hline
	\endfirsthead
	
	\continuecaption{Продолжение таблицы \ref{tab:js_stat_admin}}
	\hline
	\textbf{Название} & \textbf{Доступ} & \textbf{Описание} \\
	\hline
	\endhead
	
	\hline
	\endlastfoot
	
	handleFacultyStat & public & Запрос статистики по факультету за выбранный период, отображение результатов. \\
	\hline
	handleDepartmentStat & public & Запрос статистики по кафедре. \\
	\hline
	handleUserStat & public & Запрос статистики по пользователю. \\
	\hline
	handleCreateFaculty & public & Отправка данных для создания факультета (только администратор). \\
	\hline
	handleCreateDepartment & public & Отправка данных для создания кафедры (только администратор). \\
	\hline
\end{xltabular}

\subsubsection{Вспомогательные функции}

В таблице~\ref{tab:js_utils} описаны вспомогательные функции.

\begin{xltabular}{\textwidth}{|l|l|X|}
	\caption{Вспомогательные функции\label{tab:js_utils}}\\
	\hline
	\textbf{Название} & \textbf{Доступ} & \textbf{Описание} \\
	\hline
	\endfirsthead
	
	\continuecaption{Продолжение таблицы \ref{tab:js_utils}}
	\hline
	\textbf{Название} & \textbf{Доступ} & \textbf{Описание} \\
	\hline
	\endhead
	
	\hline
	\endlastfoot
	
	bindTextareaCounter & public & Реализация счётчика символов для текстовых полей с ограничением максимальной длины. \\
	\hline
	validateFacultyCode & internal & Проверка формата кода факультета (регулярное выражение \^\textbackslash d\{2\}\textbackslash .\textbackslash d\{2\}\textbackslash .\textbackslash d\{2\}\$). \\
	\hline
	activateSection & public & Переключение между видимыми секциями (авторизация, профиль, проверка, титульный лист, статистика, РИО, администрирование). \\
	\hline
\end{xltabular}

\subsection{Модуль images/}

Каталог images/ содержит графические ресурсы интерфейса: логотипы, иконки навигации, анимации загрузки и вспомогательные визуальные компоненты.

\subsection{Системное тестирование клиентской части}

Системное тестирование клиентской части проводилось в операционной системе Windows 10 с использованием браузеров Google Chrome, Mozilla Firefox и Yandex. Проверялись все основные сценарии использования, описанные в техническом задании.

\subsubsection{Сводная таблица тест-кейсов}

В таблице~\ref{tab:test_cases} представлена сводная таблица тест-кейсов клиентской части.

\begin{xltabular}{\textwidth}{|c|p{11.5cm}|X|}
	\caption{Сводная таблица тест-кейсов клиентской части\label{tab:test_cases}}\\
	\hline
	\textbf{ID} & \textbf{Название тест-кейса} & \textbf{Результат} \\
	\hline
	\endfirsthead
	
	\continuecaption{Продолжение таблицы \ref{tab:test_cases}}
	\hline
	\textbf{ID} & \textbf{Название тест-кейса} & \textbf{Результат} \\
	\hline
	\endhead
	
	\hline
	\endlastfoot
	
	TC-01 & Вход в систему с корректными данными & Пройден \\
	\hline
	TC-02 & Регистрация нового пользователя & Пройден \\
	\hline
	TC-03 & Загрузка и проверка корректного DOCX-документа & Пройден \\
	\hline
	TC-04 & Загрузка и проверка документа с ошибками оформления & Пройден \\
	\hline
	TC-05 & Генерация титульного листа методических указаний & Пройден \\
	\hline
	TC-06 & Генерация титульного листа учебного пособия & Пройден \\
	\hline
	TC-07 & Генерация титульного листа монографии & Пройден \\
	\hline
	TC-08 & Отправка документов в РИО (полный цикл) & Пройден \\
	\hline
	TC-09 & Получение статистики по факультету & Пройден \\
	\hline
	TC-10 & Создание факультета администратором & Пройден \\
	\hline
\end{xltabular}

\subsubsection{TC-01. Вход в систему с корректными данными}

\textbf{Предусловие:} пользователь зарегистрирован, открыта страница авторизации.

\textbf{Шаги:}
\begin{enumerate}
	\item Ввести email и пароль.
	\item Нажать кнопку «Войти».
\end{enumerate}

\textbf{Ожидаемый результат:} интерфейс переключается на страницу профиля, в локальном хранилище появляются access и refresh токены.

\textbf{Фактический результат:} пройден.

\subsubsection{TC-02. Регистрация нового пользователя}

\textbf{Предусловие:} открыта страница регистрации.

\textbf{Шаги:}
\begin{enumerate}
	\item Заполнить поля: ФИО, email, пароль, код факультета, название кафедры, выбрать роль.
	\item Нажать кнопку «Зарегистрироваться».
\end{enumerate}

\textbf{Ожидаемый результат:} появляется сообщение об успешной регистрации, пользователь перенаправляется на страницу входа.

\textbf{Фактический результат:} пройден.

\subsubsection{TC-03. Загрузка и проверка корректного DOCX-документа}

\textbf{Предусловие:} пользователь авторизован, открыта страница проверки документов.

\textbf{Шаги:}
\begin{enumerate}
	\item Выбрать тип документа «Методические указания».
	\item Загрузить файл через кнопку или перетаскиванием в область drag-and-drop.
	\item Нажать кнопку «Проверить методические указания».
\end{enumerate}

\textbf{Ожидаемый результат:} отображается анимированный индикатор загрузки, затем появляется результат «Ошибок не найдено». В разделе «Мои файлы» появляется запись о проверенном документе.

\textbf{Фактический результат:} пройден.

\subsubsection{TC-04. Загрузка и проверка документа с ошибками оформления}

\textbf{Предусловие:} пользователь авторизован, открыта страница проверки документов.

\textbf{Шаги:}
\begin{enumerate}
	\item Выбрать тип документа «Методические указания».
	\item Загрузить файл с намеренными нарушениями оформления (неверный шрифт, отсутствие отступа).
	\item Нажать кнопку «Проверить методические указания».
\end{enumerate}

\textbf{Ожидаемый результат:} система выявляет ошибки, отображает их количество, предоставляет ссылки на скачивание исправленного документа и PDF-отчёта.

\textbf{Фактический результат:} пройден.

\subsubsection{TC-05. Генерация титульного листа методических указаний}

\textbf{Предусловие:} пользователь авторизован, открыта страница генерации титульных листов.

\textbf{Шаги:}
\begin{enumerate}
	\item Выбрать тип «Методические указания».
	\item Заполнить поля: название, дисциплина, направление подготовки, составитель, рецензенты (один или несколько), УДК, описание.
	\item Нажать кнопку «Сгенерировать».
\end{enumerate}

\textbf{Ожидаемый результат:} браузер скачивает DOCX-файл с титульным листом, оформленным по шаблону (две страницы, правильное форматирование).

\textbf{Фактический результат:} пройден.

\subsubsection{TC-06. Генерация титульного листа учебного пособия}

\textbf{Предусловие:} пользователь авторизован, открыта страница генерации титульных листов.

\textbf{Шаги:}
\begin{enumerate}
	\item Выбрать тип «Учебное пособие».
	\item Заполнить поля: автор, название, город, год, рецензенты, направления подготовки, значение А, ISBN, УДК, ББК, описание.
	\item Нажать кнопку «Сгенерировать».
\end{enumerate}

\textbf{Ожидаемый результат:} браузер скачивает DOCX-файл с правильно оформленным титульным листом учебного пособия.

\textbf{Фактический результат:} пройден.

\subsubsection{TC-07. Генерация титульного листа монографии}

\textbf{Предусловие:} пользователь авторизован, открыта страница генерации титульных листов.

\textbf{Шаги:}
\begin{enumerate}
	\item Выбрать тип «Монография».
	\item Заполнить поля: авторы (один или несколько), название монографии, город, год, УДК, ББК, ISBN, описание.
	\item Нажать кнопку «Сгенерировать».
\end{enumerate}

\textbf{Ожидаемый результат:} браузер скачивает DOCX-файл с титульным листом монографии (авторы в одну строку, название, выходные данные, аннотация).

\textbf{Фактический результат:} пройден.

\subsubsection{TC-08. Отправка документов в РИО (полный цикл)}

\textbf{Предусловие:} пользователь авторизован, открыта страница «Отправить в РИО».

\textbf{Шаги:}
\begin{enumerate}
	\item На шаге 1 загрузить два любых файла (PDF, TXT, DOCX) и нажать «Загрузить файлы и продолжить».
	\item На шаге 2 загрузить основной DOCX-файл, дождаться сообщения об успешной проверке, нажать «Проверить файл и продолжить».
	\item На шаге 3 ввести комментарий (не более 1000 символов) и нажать «Отправить в РИО».
\end{enumerate}

\textbf{Ожидаемый результат:} система отображает сообщение «Файлы успешно отправлены в РИО». При проверке почтового ящика РИО обнаруживается письмо с тремя вложениями.

\textbf{Фактический результат:} пройден.

\subsubsection{TC-09. Получение статистики по факультету}

\textbf{Предусловие:} пользователь авторизован, открыта страница статистики.

\textbf{Шаги:}
\begin{enumerate}
	\item Выбрать вкладку «По факультету», ввести код факультета (например, 09.03.04).
	\item Указать даты начала и конца периода.
	\item Нажать кнопку «Посчитать».
\end{enumerate}

\textbf{Ожидаемый результат:} система отображает количество методических указаний, учебных пособий, монографий и общее число документов за выбранный период.

\textbf{Фактический результат:} пройден.

\subsubsection{TC-10. Создание факультета администратором}

\textbf{Предусловие:} пользователь с ролью администратора авторизован, открыта административная панель.

\textbf{Шаги:}
\begin{enumerate}
	\item Заполнить поля: название факультета, код (в формате 00.00.00), ФИО декана.
	\item Нажать кнопку «Создать факультет».
\end{enumerate}

\textbf{Ожидаемый результат:} система отображает сообщение об успешном создании, факультет появляется в базе данных.

\textbf{Фактический результат:} пройден.

\subsubsection{Результаты тестирования}

В ходе системного тестирования клиентской части были успешно проверены все основные сценарии использования: аутентификация, загрузка и проверка документов, генерация титульных листов, отправка в РИО, статистика и администрирование. Интерфейс корректен, динамические формы (добавление рецензентов, направлений, авторов) работают согласно спецификации, drag-and-drop загрузка файлов функционирует во всех поддерживаемых браузерах. Ошибок в работе клиентской части не выявлено. Клиентское приложение полностью соответствует требованиям технического задания и готово к промышленной эксплуатации.